\section{Complete higher-dimensional saturation geometry}\label{sec:spectral-saturation}

For a state $\rho$ on $\mathbb C^d\otimes\mathbb C^d$, write
\[
 f_d(\rho)=\max_{|e\rangle\ \text{maximally entangled}}
                  \langle e|\rho|e\rangle.
\]
Let $\alpha_1\ge\cdots\ge\alpha_d\ge0$ and
$\beta_1\ge\cdots\ge\beta_d\ge0$ be its marginal eigenvalues. Put
\[
 m_i=\min(\alpha_i,\beta_i),\qquad
 L=\sum_i m_i,\qquad \delta=1-L,
\]
\[
 I_A=\{i:\alpha_i>\beta_i\},\qquad
 I_B=\{i:\beta_i>\alpha_i\},\qquad
 |z_m\rangle=\sum_i\sqrt{m_i}|ii\rangle.
\]
The sets $I_A,I_B$ are disjoint. The sharp marginal-spectrum bound is
\begin{equation}\label{eq:marginal-fef}
 f_d(\rho)\le {1\over d}\left(\sum_i\sqrt{m_i}\right)^2.
\end{equation}

\subsection{All equality cases, including repeated and zero eigenvalues}

\paragraph{Theorem (complete saturation normal form).}
Equality in \eqref{eq:marginal-fef} holds if and only if, up to local unitaries,
\begin{equation}\label{eq:all-saturators}
 \rho=|z_m\rangle\langle z_m|+R,\qquad
 R\succeq0,\qquad
 \operatorname{supp}R\subseteq
       \operatorname{span}\{|ij\rangle:i\in I_A,\ j\in I_B\},
\end{equation}
where $\rho_A,\rho_B$ have the specified spectra $\alpha,\beta$.
There is no simplicity or full-rank assumption on either spectrum.

The remaining spectral condition has a smaller, explicit form. Write
$M_A=\operatorname{diag}(m_i:i\in I_A)$ and
$M_B=\operatorname{diag}(m_i:i\in I_B)$. Regarding $R$ as an operator on
$\mathbb C^{I_A}\otimes\mathbb C^{I_B}$, the permitted residuals are exactly
\[
 R\succeq0,\quad
 \operatorname{spec}(M_A+\operatorname{Tr}_B R)
                  =\{\alpha_i:i\in I_A\},\quad
 \operatorname{spec}(M_B+\operatorname{Tr}_A R)
                  =\{\beta_i:i\in I_B\}.
\]
In particular $\operatorname{Tr}R=\delta$. Either index set can be empty
only when $\alpha=\beta$; in that case $R=0$ and the saturator is pure.
The displayed characterization is of states with fixed marginal spectra.
For fixed marginal matrices one additionally imposes their matrix equalities.

\paragraph{Idea of the proof.} A saturating state contains the coherent
vector $z_m$ with the largest weights the marginals allow. Whatever is left
over is a positive residual whose marginals only add to the parts of the
spectra where one marginal exceeds the other. The lemma below handles
repeated eigenvalues, and positivity then confines the residual to the
rectangle $I_A\times I_B$, which the maximally entangled target does not
see.

\paragraph{A positive-perturbation lemma.}
Let $M=\operatorname{diag}(m_1,\ldots,m_d)$ with decreasing $m_i$, and
$Y\succeq0$. Suppose $M+Y$ has decreasing eigenvalues $a_i$.
For a maximal constant block $m_{k+1}=\cdots=m_{k+r}=\lambda$, let
$E_\lambda$ be its eigenspace and let
\[
 s=\#\{i\in\{k+1,\ldots,k+r\}:a_i>\lambda\}.
\]
Then $\dim(E_\lambda\cap\ker Y)\ge r-s$.

\emph{Proof.} The case $s=r$ is immediate. Otherwise monotonicity gives
$a_i\ge m_i$, so $M+Y$ has exactly $k+s$ eigenvalues strictly greater than
$\lambda$. On $V=E_{>\lambda}\oplus E_\lambda$, the compression of
$M+Y-\lambda I$ is positive semidefinite. Its kernel is exactly
$E_\lambda\cap\ker Y$: both nonnegative quadratic-form summands must
vanish, and a zero quadratic form of a positive operator is in its kernel.
Consequently this compression has a strictly positive subspace of dimension
$k+r-\dim(E_\lambda\cap\ker Y)$. The min--max principle bounds that
dimension by the number $k+s$ of positive eigenvalues of $M+Y-\lambda I$.
This proves the lemma. This argument also covers $\lambda=0$.

\paragraph{Proof of the theorem.}
Take a maximizing maximally entangled vector, and use local unitaries so it
is $|\Phi_d\rangle=d^{-1/2}\sum_i|ii\rangle$. Put
$t=\langle\Phi_d|\rho|\Phi_d\rangle$ and
$|w\rangle=\rho|\Phi_d\rangle/\sqrt t$. The zero case is irrelevant:
$L>0$ for probability spectra, so the upper bound being saturated is positive.
Positive-matrix Cauchy--Schwarz gives
\[
 R:=\rho-|w\rangle\langle w|\succeq0,\qquad
 \langle\Phi_d|w\rangle=\sqrt t.
\]
If $W$ is the coefficient matrix of $w$ and $q_i$ are its squared singular
values in decreasing order, partial trace and eigenvalue monotonicity give
$q_i\le m_i$. Therefore
\[
 \sqrt{dt}=\operatorname{Re}\operatorname{Tr}W
 \le\operatorname{Tr}\sqrt{W^\dagger W}
 =\sum_i\sqrt{q_i}\le\sum_i\sqrt{m_i}.
\]
Equality forces $q_i=m_i$ for every $i$ and $W\succeq0$.
For the latter assertion, write $W=UH$ in polar form: equality in
$\operatorname{Re}\operatorname{Tr}UH\le\operatorname{Tr}H$ forces
$U$ to be the identity on the range of $H$. A transformation
$U_0\otimes\overline{U_0}$ preserves $\Phi_d$ and diagonalizes $W$.
Thus $w=z_m$, and
\[
 \rho_A=M+R_A,\qquad \rho_B^T=M+R_B^T,
\]
with both residual marginals positive.

Consider any maximal constant block of $m$. The indices in that block with
$\alpha_i>m_i$ form a prefix, as do those with $\beta_i>m_i$. These two
prefixes cannot both be nonempty, since that would contradict
$m_i=\min(\alpha_i,\beta_i)$ at the first position of the block. Hence
at most one residual marginal is active on any such block. Apply the lemma
to both residual marginals. The inactive marginal annihilates the entire
block. The active marginal has at least as many kernel vectors in that
block as there are positions without surplus. Choose an orthonormal basis
of the block whose last vectors lie in that kernel. This basis choice
preserves both $M$ and $\Phi_d$, so it can be made independently on every
constant block. In the resulting bases
\[
 \operatorname{supp}R_A\subseteq\mathbb C^{I_A},\qquad
 \operatorname{supp}R_B\subseteq\mathbb C^{I_B}.
\]
Positivity upgrades marginal support to joint support. Indeed, for any local
kernel vector $u$, the nonnegative sum
$\sum_j\langle u,j|R|u,j\rangle=\langle u|R_A|u\rangle=0$
forces every $|u,j\rangle$ into $\ker R$. Applying this on both sides
gives the rectangle in \eqref{eq:all-saturators}.

Conversely, a positive residual on that rectangle is orthogonal to every
$|ii\rangle$, because $I_A\cap I_B=\varnothing$. Therefore
\[
 \langle\Phi_d|\rho|\Phi_d\rangle
       ={1\over d}\left(\sum_i\sqrt{m_i}\right)^2,
\]
and the spectral upper bound makes this the fully entangled fraction.
The smaller spectral conditions follow by splitting each marginal into
its active residual block and the unaffected diagonal complement.

\subsection{The same equality set for every full-Schmidt-rank target orbit}

Let $p_1\ge\cdots\ge p_d\ge0$ be a probability vector and define
\[
 F_p(\rho)=\max_{\substack{|x\rangle:\ \text{squared Schmidt coefficients }p}}
                  \langle x|\rho|x\rangle.
\]
The sharp overlap bound is
\[
 F_p(\rho)\le\left(\sum_i\sqrt{p_i m_i}\right)^2.
\]
For fixed marginal spectra $\alpha,\beta$, the following are equivalent:
\begin{enumerate}
\item $\lambda_{\max}(\rho)=L=\sum_i m_i$;
\item $\rho$ has the complete normal form \eqref{eq:all-saturators};
\item $f_d(\rho)=d^{-1}(\sum_i\sqrt{m_i})^2$;
\item for at least one decreasing probability vector $p$ with all $p_i>0$,
      $F_p(\rho)=(\sum_i\sqrt{p_i m_i})^2$;
\item for every decreasing probability vector $p$, including those with zero entries,
      $F_p(\rho)=(\sum_i\sqrt{p_i m_i})^2$.
\end{enumerate}
Thus the complete maximizing set is independent of every strictly positive
target Schmidt profile. This identifies all maximizers of the previously
known maximal-eigenvalue bound, not merely one family attaining its value.

\emph{Proof.} Suppose first that $\lambda_{\max}(\rho)=L$ and choose a
normalized top eigenvector $v$. The vector $w=\sqrt L\,v$ satisfies
$\rho\succeq|w\rangle\langle w|$. If $q_i$ are its decreasing squared
Schmidt coefficients, marginal positivity gives $q_i\le m_i$. But
$\sum_i q_i=L=\sum_i m_i$, so every inequality is an equality. Apply
independent local unitaries to put $w=z_m$ and use the positive-perturbation
lemma above. Its block argument depends only on the marginal spectra and
$q=m$, and therefore gives \eqref{eq:all-saturators}. Conversely, a state
in that normal form has $\rho z_m=Lz_m$, since the residual rectangle is
orthogonal to $z_m$. The general bound $\lambda_{\max}(\rho)\le L$
then makes its largest eigenvalue exactly $L$.

For the target-orbit assertion, choose a maximizing target $x$ and extract
$w=\rho x/\sqrt{\langle x|\rho|x\rangle}$ as before. The singular-value
trace inequality gives
\[
 \sqrt{\langle x|\rho|x\rangle}
 \le\sum_i\sqrt{p_i q_i}\le\sum_i\sqrt{p_i m_i}.
\]
If all $p_i>0$, equality forces every $q_i=m_i$. The same block argument
again gives the normal form. Conversely, in the normal-form bases the
target $x_p=\sum_i\sqrt{p_i}|ii\rangle$ is orthogonal to $R$ and has
overlap $(\sum_i\sqrt{p_i m_i})^2$. This works also when some $p_i=0$.
The maximally entangled target is the special case $p_i=1/d$, completing
all implications.

The condition that the one target profile in item 4 have no zero entry is
needed only on the positive cap support: the same equivalence holds whenever
$p_i>0$ for every $i$ with $m_i>0$. Indeed $q_i\le m_i=0$ already fixes
all remaining extracted Schmidt coefficients. For decreasing profiles this
condition says that the target Schmidt rank is at least
$r=\min(\operatorname{rank}\rho_A,\operatorname{rank}\rho_B)$.
Some positivity condition on the active support is necessary. For $d=2$, the state
\[
 \rho={1\over2}|00\rangle\langle00|
       +{1\over2}|11\rangle\langle11|
\]
has marginal spectra $\alpha=\beta=(1/2,1/2)$ and attains the product-target
cap $F_{(1,0)}=1/2$. Nevertheless its fully entangled fraction and largest
eigenvalue are both $1/2$, strictly below the corresponding caps $1$.

\paragraph{Prior-art boundary for the equivalence.}
Klyachko~\cite{klyachko}, Theorem 6.3.1, already proved that the
maximum possible largest eigenvalue for prescribed marginal spectra is
$\sum_i\min(\alpha_i,\beta_i)$. Immediately after that theorem he raised
the problem of describing its maximizers and discussed a representation-theoretic
recurrence for the associated spectra. The complete positive-residual
normal form and its equivalence with every full-Schmidt-rank overlap
maximization are the additional results proved here. A bounded search for
subsequent equality classifications found no matching statement; that search
does not establish worldwide priority or certify that the classification
problem remained open until this work.

\subsection{Exact rank-one coupling capacity in every dimension}

For positive integers $n\le m$, define the uniform-marginal coupling capacity
of a normalized pure target $x\in\mathbb C^n\otimes\mathbb C^m$ by
\[
 \mathsf a_{n,m}(x)=\max_{\substack{\rho\succeq0\\
                     \rho_A=I_n/n,\ \rho_B=I_m/m}}
                   \langle x|\rho|x\rangle.
\]
If $s_1,\ldots,s_n\ge0$ are the Schmidt amplitudes of $x$, then
\begin{equation}\label{eq:uniform-capacity}
 \boxed{\mathsf a_{n,m}(x)={1\over m}\left(\sum_{i=1}^n s_i\right)^2.}
\end{equation}
Equivalently, for arbitrary local dimensions the denominator is
$\max(n,m)$ and the sum runs over the $\min(n,m)$ Schmidt amplitudes.
For $n=m=d$ and every $s_i>0$, the optimizing coupling is unique:
it is the maximally entangled pure state in the target's Schmidt bases.

\emph{Elementary proof.} Use local unitaries to write
$x=\sum_{i=1}^n s_i|ii\rangle$. Put $d_i=\rho_{ii,ii}$.
The Bob marginal gives $0\le d_i\le1/m$, while positivity gives
$\operatorname{Re}\rho_{ii,jj}\le\sqrt{d_i d_j}\le1/m$.
Consequently
\[
 \langle x|\rho|x\rangle
   \le\left(\sum_i s_i\sqrt{d_i}\right)^2
   \le {1\over m}\left(\sum_i s_i\right)^2.
\]
This proof also permits zero Schmidt amplitudes. For sharpness, let
\[
 z={1\over\sqrt m}\sum_{i=1}^n|ii\rangle,\qquad
 \rho_*=|z\rangle\langle z|+
        {1\over nm}\sum_{i=1}^n\sum_{j=n+1}^m|ij\rangle\langle ij|.
\]
The sum is empty if $n=m$. Its Alice marginal is
$I_n/m+(m-n)I_n/(nm)=I_n/n$; its Bob marginal is $I_m/m$.
It is positive and its target overlap is the right side of
\eqref{eq:uniform-capacity}. In the square case with every $s_i>0$,
equality forces every $d_i=1/d$. These $d$ populations already sum to one,
so all other populations, and hence their rows and columns by positivity,
vanish. Equality in the off-diagonal estimates forces the remaining block
to have every entry $1/d$, giving the unique pure Bell coupling.

More generally, if every $s_i>0$ and $n<m$, the entire rectangular
optimizer face is
\[
 \rho=|z\rangle\langle z|+R,\quad R\succeq0,\quad
 \operatorname{supp}R\subseteq\mathbb C^n\otimes
                    \operatorname{span}\{|n+1\rangle,\ldots,|m\rangle\},
\]
\[
 R_A=(1/n-1/m)I_n,\qquad R_B=P_{\mathrm{tail}}/m.
\]
Indeed equality forces $\rho_{ii,ii}=1/m$ for all $i\le n$ and the
correlated principal block $zz^*$. The fixed Bob marginal forces every
other population in the first $n$ Bob columns to vanish. Positivity leaves
only a possible coherent cross vector $v$ between $z$ and the tail, but
the Bob off-diagonal entries between the first $n$ columns and the tail
are $m^{-1/2}v_{ij}$ or their conjugates; the prescribed diagonal Bob
marginal therefore forces $v=0$. The remaining partial traces are exactly
the displayed ones. Conversely every such residual attains the bound.
The face has affine dimension
\[
 (n^2-1)((m-n)^2-1).
\]
For $n<m$ the optimizer is unique precisely when $n=1$ or $m=n+1$;
for $n=m$ it is the unique pure Bell state. The two-by-three attainer below
is consequently unique for its displayed target.

\paragraph{Published capacity comparison.}
Proposition 3.5 of Skalski et al.~\cite{capacity2025} proves the lower bound
corresponding to \eqref{eq:uniform-capacity} in the square case and proves
equality when $d=2$. The formula above establishes equality for every $d$.
The paper does not label this particular all-dimensional extension as an
open problem. Its capacity $\alpha(E_x)$ equals $\mathsf a_{n,m}(x)$:
the paper's normalized-trace density $A_\sigma$ corresponds to the ordinary
trace-one matrix $\rho=A_\sigma/(nm)$, and the normalized-trace marginal
functionals correspond exactly to $I_n/n,I_m/m$.

The same paper's printed Theorem 3.7(ii), stated for $n\le m$, says that a
maximally entangled target has capacity one. With the paper's definition
of rectangular maximal entanglement, its Schmidt amplitudes are all
$1/\sqrt n$. Formula \eqref{eq:uniform-capacity} gives the corrected value
$n/m$, so the printed statement requires either $n=m$ or replacement of
one by $n/m$. In fact Cauchy--Schwarz in \eqref{eq:uniform-capacity}
proves the exact rectangular characterization
\[
 x\text{ is maximally entangled in }\mathbb C^n\otimes\mathbb C^m
 \quad\Longleftrightarrow\quad \mathsf a_{n,m}(x)=n/m.
\]

For a concrete counterexample to the printed rectangular value, take
\[
 x={|00\rangle+|11\rangle\over\sqrt2}\in\mathbb C^2\otimes\mathbb C^3,
 \qquad
 \rho_*={2\over3}|x\rangle\langle x|
       +{1\over6}\bigl(|02\rangle\langle02|+|12\rangle\langle12|\bigr).
\]
Both marginals are uniform, and the elementary bound above proves
$\mathsf a_{2,3}(x)=2/3<1$ exactly. The printed statement has this form in arXiv:2211.02401v3; no erratum was found in a bounded search. This is a correction of the displayed rectangular
statement; it does not affect the paper's square-dimensional statement
or its general duality theorem.

For completeness, the same paper's projective coupling capacity has the
particularly simple exact expression
\[
 \gamma(E_x)={r\over m},\qquad r=\operatorname{Schmidt\ rank}(x),\quad n\le m.
\]
Its defining feasibility condition is
$(P^\perp\otimes Q^\perp)x=0$ for local orthogonal projections $P,Q$.
If $X$ is the coefficient matrix of $x$, this condition writes $X$ as a
sum of a matrix with range in $\operatorname{ran}P$ and one whose right
support has dimension $\operatorname{rank}Q$. Thus
$r\le\operatorname{rank}P+\operatorname{rank}Q$, and the normalized cost
satisfies
\[
 {\operatorname{rank}P\over n}+{\operatorname{rank}Q\over m}
 \ge {r\over m}.
\]
Taking $P=0$ and $Q$ equal to the Bob Schmidt-support projection attains
this value. Hence $\alpha(E_x)=\gamma(E_x)$ exactly when all nonzero
Schmidt amplitudes are equal, by Cauchy--Schwarz.

The paper ends by asking whether its Proposition 3.10 extends beyond
dimension two. Part (i), which would say that a subspace supporting a
uniform-marginal coupling must contain a maximally entangled vector, has
a negative answer already in dimension three. Let $F$ be the swap and
$P_-=(I-F)/2$ the projection onto the antisymmetric subspace. The state
\[
 \rho_-={P_-\over3}={I-F\over6}
\]
has both marginals $I_3/3$ and is supported on $P_-$, so
$\alpha(P_-)=1$. Every vector in that subspace has a skew-symmetric
$3\times3$ coefficient matrix, whose determinant is zero since
$\det C=\det(-C^T)=-\det C$. It therefore has Schmidt rank at most two
and cannot be maximally entangled in dimension three. This uses the
standard antisymmetric state as an exact counterexample to the proposed
extension; it is not a new construction of that state. The formula for
$\gamma$ likewise gives the precise higher-dimensional replacement for
the rank-one part (ii). These observations are immediate consequences of
the displayed definitions and do not require a further conjectural step.

\subsection{The aligned face and sharp rank restriction}
For decreasing diagonal marginal matrices and a fixed target with positive
diagonal Schmidt amplitudes, the complete face is the coherent core plus
all residual couplings with marginals $\operatorname{diag}(\alpha-m)$ and
$\operatorname{diag}(\beta-m)$. This is the ordered specialization of
Theorems~\ref{thm:optimizer-face} and~\ref{thm:boundary-face}. Even when
some $m_i=0$, at most one marginal has a surplus in the final zero block,
so the possible zero-block cross rectangle in the general unordered
boundary theorem is empty.

For every spectral saturator the normal form gives
\[
 \operatorname{rank}\rho=1+\operatorname{rank}R
 \le1+|I_A||I_B|\le1+\lfloor d^2/4\rfloor.
\]
This is sharp for every $d\ge2$: choose a positive strictly decreasing
$m$ of sum $1-\delta$, partition the surplus indices as evenly as possible,
and distribute the two residual masses positively on opposite sides.
For sufficiently small $\delta$, both marginal lists remain decreasing.
The product residual is positive definite on its rectangle and attains
the bound. This is consistent with the more general fixed-target theorem
in Section~\ref{sec:all-optimizer-rank}; it describes every saturator,
whereas Theorem~\ref{thm:low-rank-pure-target} guarantees a smaller-rank
attaining representative.

\subsection{Every saturator has the largest possible Schmidt number}

\paragraph{Theorem.}
Every state saturating \eqref{eq:marginal-fef} has
\[
 \operatorname{SN}(\rho)
       =\min\{\operatorname{rank}\rho_A,\operatorname{rank}\rho_B\}.
\]
This conclusion includes all degeneracies and rank-deficient marginal states.

\emph{Proof.} Let $r=\#\{i:m_i>0\}$, which equals the displayed minimum
of local ranks. In the normal form, the support of $\rho$ is the orthogonal
sum of $\operatorname{span}\{z_m\}$ and $\operatorname{supp}R$.
Every vector in that support has a coefficient matrix
\[
 C=a\operatorname{diag}(\sqrt{m_1},\ldots,\sqrt{m_d})+N,
 \qquad N_{ij}=0\quad\hbox{unless }i\in I_A, j\in I_B.
\]
Reorder the first $r$ indices by placing $I_A$ first and $I_B$ last.
The leading $r\times r$ block becomes upper triangular apart from its
diagonal, and therefore
\[
 \det C_{[r],[r]}=a^r\prod_{i=1}^r\sqrt{m_i}.
\]
For $a\ne0$ this minor is nonzero, so the vector has Schmidt rank at least
$r$. Every pure-state decomposition of $\rho$ contains a vector with
$a\ne0$, because its expectation on $z_m$ is nonzero. Hence its maximum
Schmidt rank is at least $r$, proving $\operatorname{SN}(\rho)\ge r$.
Every vector in the support of any bipartite state lies in the tensor product
of its marginal supports, giving the reverse bound $\operatorname{SN}(\rho)
\le\min(\operatorname{rank}\rho_A,\operatorname{rank}\rho_B)=r$.

\subsection{Fixed partial-transpose inertia and one-copy distillability}

Write $\Gamma$ for the partial transpose on Bob. If every $m_i>0$, every
saturator has the exact partial-transpose inertia
\[
 \boxed{\quad
 n_+(\rho^\Gamma)={d(d+1)\over2},\qquad
 n_-(\rho^\Gamma)={d(d-1)\over2},\qquad n_0(\rho^\Gamma)=0.
 \quad}
\]
The statement permits an arbitrary positive quantum residual in the normal
form. If exactly $r$ cap eigenvalues are positive, then
$n_-(\rho^\Gamma)\ge r(r-1)/2$. Whenever $r\ge2$, a single copy can be
locally projected with positive probability to an entangled two-qubit state;
thus every such saturator is one-copy distillable.

\emph{Proof.} Put $T=I_A\times I_B$ and let
$T'=\{(j,i):(i,j)\in T\}$. These coordinate sets are disjoint. In the
paired bases of $T$ and $T'$, the corresponding partial-transpose block is
\[
 H=\begin{pmatrix}C&D\\D&0\end{pmatrix},\qquad
 C=(R^\Gamma)_{T,T}=C^\dagger,\qquad
 D=\operatorname{diag}(\sqrt{m_i m_j}:(i,j)\in T)>0.
\]
There are no couplings from this block to its complement. The invertible
shear
\[
 S=\begin{pmatrix}I&0\\-\tfrac12D^{-1}C&I\end{pmatrix}
 \quad\hbox{satisfies}\quad
 S^\dagger H S=\begin{pmatrix}0&D\\D&0\end{pmatrix}.
\]
Sylvester's law of inertia therefore gives $|T|$ positive and $|T|$
negative eigenvalues on this block. Every remaining unordered pair
$\{i,j\}$ gives the unchanged two-dimensional block with eigenvalues
$\pm\sqrt{m_i m_j}$, and each diagonal vector $|ii\rangle$ gives the
positive eigenvalue $m_i$. Counting these blocks proves the full-rank
formula. Applying the same argument to the principal compression onto the
first $r$ indices on both sides, followed by eigenvalue interlacing, proves
the lower bound when some $m_i$ vanish.

Finally fix any $i\ne j$ with $m_i,m_j>0$ and project both local spaces
onto $\operatorname{span}\{|i\rangle,|j\rangle\}$. The coherent term
survives. The residual rectangle contributes at most one of the two
populations $|ij\rangle\langle ij|$ and $|ji\rangle\langle ji|$.
The projected state's partial transpose consequently has a principal
two-dimensional block of the form
\[
 \begin{pmatrix}\kappa&\sqrt{m_i m_j}\\
                 \sqrt{m_i m_j}&0\end{pmatrix}
 \quad\hbox{or its reversed ordering},\qquad \kappa\ge0,
\]
whose determinant is $-m_i m_j<0$. Its negative direction has Schmidt
rank two. The local projection succeeds with probability at least
$m_i+m_j>0$, proving the asserted one-copy distillability.

\subsection{An exact convex-roof entanglement value}

For full-rank $m$, use the standard normalization of the $G$-concurrence:
for a normalized pure state with coefficient matrix $C$,
$G_d(C)=d|\det C|^{2/d}$. Let $G_d(\rho)$ be its convex roof and
$G_d^{\mathrm a}(\rho)$ the concave roof over pure-state decompositions.
Then every saturator has the exact values
\[
 G_d(\rho)=G_d^{\mathrm a}(\rho)
                   =d\left(\prod_{i=1}^d m_i\right)^{1/d}.
\]
In fact every pure-state decomposition has the same average $G$-concurrence.
For any normalized support vector $C=a\operatorname{diag}\sqrt m+N$,
the triangular determinant identity gives
$G_d(C)=d|a|^2(\prod_i m_i)^{1/d}$. For a decomposition of $\rho$,
the expectation of the projection onto $z_m/\sqrt L$ gives
$\sum_k p_k|a_k|^2=1$. Averaging proves the formula without optimization.

This affine-roof mechanism has prior art: Gour and Guo~\cite{gourguo}, Theorem 7, characterize equality of $G$-concurrence and assistance using
nilpotent subspaces. The rectangle matrices here square to zero after
diagonal rescaling. The new connection established above is that every
marginal-spectrum FEF maximizer lies in this class, with the displayed
entanglement value fixed entirely by the two marginal spectra.

\subsection{Explicit quasidistillation, including quantum residuals}

Let $r=\#\{i:m_i>0\}$, put $\mu=\min_{i\le r}m_i$, and write a saturator in
the normal form. Let $D$ be diagonal with $D_i=m_i^{-1/4}$ for $i\le r$
and $D_i=0$ otherwise, and put
$\widetilde R=(D\otimes D)R(D\otimes D)$, with
$r_R=\operatorname{Tr}\widetilde R$. For $0<\varepsilon\le1$, define
diagonal filters
\[
 A_\varepsilon=\mu^{1/4}D\operatorname{diag}(\varepsilon^{a_i}),\qquad
 B_\varepsilon=\mu^{1/4}D\operatorname{diag}(\varepsilon^{1-a_i}),
\]
where $a_i=1$ on $I_A$, $a_i=0$ on $I_B$, and $a_i=1/2$ otherwise.
Both are contractions, since every retained diagonal factor
$\mu^{1/4}m_i^{-1/4}\varepsilon^{a_i}$ is at most one and the other
entries vanish. Direct multiplication
gives the unnormalized output and its success probability:
\[
 (A_\varepsilon\otimes B_\varepsilon)\rho
 (A_\varepsilon\otimes B_\varepsilon)^\dagger
 =\mu\varepsilon^2\bigl(r|\Phi_r\rangle\langle\Phi_r|
                               +\varepsilon^2\widetilde R\bigr),
\]
\[
 P_\varepsilon=\mu\varepsilon^2(r+\varepsilon^2r_R),\qquad
 F_\varepsilon:=\langle\Phi_r|\rho_\varepsilon|\Phi_r\rangle
                         ={r\over r+\varepsilon^2r_R}.
\]
Here $\Phi_r=r^{-1/2}\sum_{i=1}^r|ii\rangle$ is embedded in the original
space. Because $\widetilde R\perp\Phi_r$, the trace distance is exact:
\[
 {1\over2}\|\rho_\varepsilon-|\Phi_r\rangle\langle\Phi_r|\|_1
                =1-F_\varepsilon.
\]
For $\varepsilon^2\lambda_{\max}(\widetilde R)\le r$, the displayed Bell
vector is also a top eigenvector, and the fully entangled fraction within
the surviving $r\times r$ subspace equals $F_\varepsilon$.
Thus every saturator can approach the maximally entangled state of the
largest dimension permitted by its marginal ranks, even when its
residual $R$ is entangled and has off-diagonal coherences. For $R\ne0$ the
exact protocol relation, for sufficiently large $F$, is
\[
 P={\mu r^2\over r_R}{1-F\over F^2}\qquad(r_R>0).
\]
This is an attained protocol curve, not a claim of globally optimal success
probability. If $r_R=0$, the projected output is already exactly $\Phi_r$
with positive success probability $\mu r\varepsilon^2$; this can happen
for a mixed original state when $r<d$. When $r=d$, a mixed saturator
cannot yield an exact $d$-dimensional Bell
state in a single nonzero-probability local-filter branch: full Schmidt
rank of the output would require both local filters to be invertible,
whereas invertible filtering preserves the state's rank.

\subsection{Nonuniqueness even for fixed, simple marginal matrices}

Take $d=4$ and
\[
 \alpha=(2/5,3/10,1/5,1/10),\quad
 \beta=(7/20,33/100,17/100,3/20),\quad
 m=(7/20,3/10,17/100,1/10).
\]
All three lists have strictly decreasing positive entries. Here
$I_A=\{1,3\}$, $I_B=\{2,4\}$, $L=23/25$, and $\delta=2/25$.
Let
\[
 C={1\over20}|1\rangle\langle1|+{3\over100}|3\rangle\langle3|,
 \qquad
 D={3\over100}|2\rangle\langle2|+{1\over20}|4\rangle\langle4|.
\]
The product residual $R_0=C\otimes D/\delta$ and the pure residual
\[
 R_1=|r\rangle\langle r|,\qquad
 |r\rangle=\sqrt{1/20}|1,4\rangle+\sqrt{3/100}|3,2\rangle
\]
have exactly the same partial traces $C,D$. Consequently
\[
 \rho_s=|z_m\rangle\langle z_m|+(1-s)R_0+sR_1,\qquad0\le s\le1,
\]
all have the same fixed marginal matrices $\operatorname{diag}\alpha$ and
$\operatorname{diag}\beta$, and all attain the same sharp fully entangled
fraction. Nevertheless their purity varies:
\[
 \operatorname{Tr}\rho_s^2
 = {135713\over160000}+{3\over16000}s+{141\over32000}s^2.
\]
In particular $\rho_0$ and $\rho_1$ have purities $135713/160000$ and
$533/625$, and have ranks five and two respectively. They cannot be locally
unitarily equivalent. This gives an explicit continuum of inequivalent
saturators even with simple spectra and fixed marginal matrices. All of
them have Schmidt number four by the preceding theorem.

\subsection{Maximal Schmidt number with a small fully entangled fraction}

The Schmidt-number conclusion is stronger than certification by a large
fully entangled fraction alone. Consider the four-dimensional state
\[
 |z\rangle={1\over2}|00\rangle+
       {1\over10}(|11\rangle+|22\rangle+|33\rangle),\qquad
 \rho=|z\rangle\langle z|+{6\over25}\sum_{j=1}^3|0j\rangle\langle0j|.
\]
Its marginals are
\[
 \rho_A=\operatorname{diag}(97/100,1/100,1/100,1/100),\qquad
 \rho_B=I/4.
\]
It attains the spectral bound with
\[
 f_4(\rho)={4\over25}<{1\over4},\qquad
 \operatorname{SN}(\rho)=4,\qquad
 G_4(\rho)=G_4^{\mathrm a}(\rho)={\sqrt5\over25}.
\]
Thus even a saturator below the usual unfiltered teleportation threshold
can possess entanglement of the maximum available dimension. The explicit
filters above have $\mu=1/100$ and $r_R=72/5$, so
\[
 F_\varepsilon={5\over5+18\varepsilon^2},\qquad
 P_\varepsilon={\varepsilon^2(25+90\varepsilon^2)\over625}.
\]
The Bell overlap equals the output fully entangled fraction for
$\varepsilon^2\le5/6$, and tends to one as the success probability tends
to zero. The state has rank four, with eigenvalues
$7/25,6/25,6/25,6/25$ and twelve zeros.

\subsection{Prior-art and formal-verification scope}

Klyachko's 2004 maximal spectral-norm theorem gives $\lambda_{\max}\rho\le L$
with sharp value and a representation-theoretic discussion of maximizers;
that value is not a new result here.
Baio et al.~\cite{baio}, Theorem 2, establish quasidistillability for a
rank-one maximally correlated component plus a \emph{classical diagonal}
residual precisely when the residual's directed support is acyclic. That
older filtering mechanism covers the classical instances above. The theorem
here classifies every saturator of the weighted marginal FEF bound, permits
arbitrary quantum residuals on the forced rectangle, and yields the exact
Schmidt-number consequence. No claim of worldwide priority is certified by
the bounded literature search.

The saturation classification, its degeneracy argument, the Schmidt-number,
partial-transpose inertia, one-copy distillability and quasidistillation
consequences, the maximizing-face dimension, and the
$G$-concurrence calculation are exact analytic proofs. They are not
yet end-to-end Lean theorems. The \texttt{PopulationOverlap} section of \texttt{AllProofs.lean}
formalizes an end-to-end actual-complex-matrix optimization using target-basis
marginal populations, with a genuine positive trace-one attaining matrix.
That result needs no spectral or singular-value hypotheses.
